% ============================================================
% style/colors.tex
%
% Цветовые палитры и смысловые алиасы темы.
%
% Идея файла:
%   1. Сначала задаются исходные палитры цветов.
%   2. Затем вводятся смысловые алиасы r..., которые используются
%      в остальных стилевых файлах.
%
% Важно:
%   p1--p9, red1--red12, myblue и т.д. — это базовые цвета.
%   rbg, rtext, rframe, rc1--rc20 — это цвета текущего оформления.
% ============================================================


% ------------------------------------------------------------
% 1. Physics Updated: синяя палитра p1--p9
% ------------------------------------------------------------

\definecolor{p1}{HTML}{caf0f8}
\definecolor{p2}{HTML}{ade8f4}
\definecolor{p3}{HTML}{90e0ef}
\definecolor{p4}{HTML}{48cae4}
\definecolor{p5}{HTML}{00b4d8}
\definecolor{p6}{HTML}{0096c7}
\definecolor{p7}{HTML}{0077b6}
\definecolor{p8}{HTML}{023e8a}
\definecolor{p9}{HTML}{03045e}


% ------------------------------------------------------------
% 2. Основные цвета страницы
% ------------------------------------------------------------
%
% pag    — главный тёмный фон страницы.
% pagtwo — дополнительный тёмный оттенок для второстепенных элементов.

\definecolor{pag}{HTML}{293133}
\definecolor{pagtwo}{HTML}{3E4547}


% ------------------------------------------------------------
% 3. Material-like цвета из Physics Updated
% ------------------------------------------------------------

\definecolor{myred}{HTML}{f44336}
\definecolor{mypink}{HTML}{e81e63}
\definecolor{mypurple}{HTML}{9c27b0}
\definecolor{mydeeppurple}{HTML}{673ab7}
\definecolor{myindigo}{HTML}{3f51b5}
\definecolor{myblue}{HTML}{2196f3}
\definecolor{mylightblue}{HTML}{03a9f4}
\definecolor{mycyan}{HTML}{00bcd4}
\definecolor{myteal}{HTML}{009688}
\definecolor{mygreen}{HTML}{4caf50}
\definecolor{mylightgreen}{HTML}{8bc34a}
\definecolor{mylime}{HTML}{cddc39}
\definecolor{myyellow}{HTML}{ffeb3b}
\definecolor{myamber}{HTML}{ffc107}
\definecolor{myorange}{HTML}{ff9800}
\definecolor{mydeeporange}{HTML}{ff5722}
\definecolor{mybrown}{HTML}{795548}
\definecolor{mygray}{HTML}{9e9e9e}
\definecolor{mybluegray}{HTML}{607d8b}


% ------------------------------------------------------------
% 4. ChemNotes2: красная палитра red1--red12
% ------------------------------------------------------------

\definecolor{red1}{HTML}{ffccd5}
\definecolor{red2}{HTML}{ffb3c1}
\definecolor{red3}{HTML}{ff8fa3}
\definecolor{red4}{HTML}{ff758f}
\definecolor{red5}{HTML}{ff4d6d}
\definecolor{red6}{HTML}{c9184a}
\definecolor{red7}{HTML}{a4133c}
\definecolor{red8}{HTML}{800f2f}
\definecolor{red9}{HTML}{590d22}
\definecolor{red10}{HTML}{c9181e}
\definecolor{red11}{HTML}{c91876}
\definecolor{red12}{HTML}{e5235a}


% ------------------------------------------------------------
% 5. Calc mindmap: оранжево-синяя палитра
% ------------------------------------------------------------

\definecolor{calcRed1}{HTML}{ff1200}
\definecolor{calcRed2}{HTML}{ff4f00}
\definecolor{calcRed3}{HTML}{ff7600}
\definecolor{calcRed4}{HTML}{ffab00}
\definecolor{calcRed5}{HTML}{ffca00}
\definecolor{calcRed6}{HTML}{ffff00}
\definecolor{calcRed7}{HTML}{bdbd42}
\definecolor{calcRed8}{HTML}{9e9e61}
\definecolor{calcRed9}{HTML}{676798}
\definecolor{calcRed10}{HTML}{3c3cc3}
\definecolor{calcRed11}{HTML}{0d0ee1}


% ------------------------------------------------------------
% 6. Цвета, сгенерированные в ChemNotes2
% ------------------------------------------------------------
%
% Эти имена выглядят технически, но оставлены без переименования,
% чтобы не потерять связь с исходной палитрой.

\definecolor{cbada55}{RGB}{186,218,85}
\definecolor{c2c2c2c}{RGB}{44,44,44}
\definecolor{cc9184a}{RGB}{201,24,74}
\definecolor{c8d354e}{RGB}{141,53,78}
\definecolor{cff4d6d}{RGB}{255,77,109}
\definecolor{cff758f}{RGB}{255,117,143}


% ============================================================
% 7. Основные смысловые алиасы темы
% ============================================================
%
% Эти цвета лучше использовать в остальных стилевых файлах.
%
% Например:
%   rbg          — основной фон;
%   rtext        — основной светлый текст;
%   rframe       — цвет декоративных рамок;
%   raccentblue  — основной синий акцент;
%   raccentred   — основной красный акцент.
%
% Если потом захочется поменять оформление, достаточно будет заменить
% цвета здесь, а не искать конкретные p4, red5 и pag по всему проекту.

\colorlet{rbg}{pag}
\colorlet{rbgsecondary}{pagtwo}

\colorlet{rtext}{white}
\colorlet{rtextmuted}{white!75}

\colorlet{rframe}{white}
\colorlet{rframesoft}{white!45}

\colorlet{raccentblue}{p4}
\colorlet{raccentbluesoft}{p3}
\colorlet{raccentbluedeep}{p7}

\colorlet{raccentred}{red5}
\colorlet{raccentredsoft}{red3}
\colorlet{raccentreddeep}{red6}

\colorlet{rgrid}{pag!95}
\colorlet{raxis}{pag!80}


% ============================================================
% 8. Быстрые цвета rc1--rc20
% ============================================================
%
% Эти алиасы нужны для быстрых выделений, графиков, декоративных
% элементов и временных цветовых экспериментов.
%
% rc1--rc10  — красная палитра.
% rc11--rc20 — синяя / голубая / фиолетовая палитра.

% ------------------------------------------------------------
% rc1--rc10: красная палитра
% ------------------------------------------------------------

\colorlet{rc1}{red1}
\colorlet{rc2}{red2}
\colorlet{rc3}{red3}
\colorlet{rc4}{red4}
\colorlet{rc5}{red5}
\colorlet{rc6}{red6}
\colorlet{rc7}{red7}
\colorlet{rc8}{red8}
\colorlet{rc9}{red9}
\colorlet{rc10}{red10}


% ------------------------------------------------------------
% rc11--rc20: синяя / голубая / фиолетовая палитра
% ------------------------------------------------------------

\colorlet{rc11}{p1}
\colorlet{rc12}{p2}
\colorlet{rc13}{p3}
\colorlet{rc14}{p4}
\colorlet{rc15}{p5}
\colorlet{rc16}{p6}
\colorlet{rc17}{p7}
\colorlet{rc18}{p8}
\colorlet{rc19}{p9}

% Добивочный фиолетовый цвет, чтобы палитра дошла до 20 цветов.
\definecolor{rc20}{HTML}{7B2CBF}